\subsection{Proper and free actions of discrete groups}

THEOREM 1. --- Let G be a discrete group acting continuously on the right on a topological space E. Suppose that every point of E has a neighborhood U such that U ∩ (U · g) = ∅ for every element g of G other than the identity element. Then the canonical map p: E → E/G makes E a principal covering of E/G with group G.

The action of G on E is free by hypothesis, so it suffices to prove that the map p is étale (I, p. 98, example 1). Let x be a point of E and let U be an open neighborhood of x such that \(U \cap U \cdot g = \varnothing\) for every element g of G distinct from the identity element. The map p is open (TG, III, p. 10, lemma 2) and induces an injective, continuous, and open map from U onto \(p(\mathrm{U})\) , hence a homeomorphism, which proves that the map p is étale.

Examples. --- 1) Let n be an integer \(\geqslant 0\) , let \(\mathbf{P}_{n}(\mathbf{R})\) be the n-dimensional projective space (TG, VI, p. 13) and let \(S_{n}\) be the unit sphere in \(R^{n+1}\) (TG, VI, p. 9). The canonical map from \(S_{n}\) onto \(\mathbf{P}_{n}(\mathbf{R})\) makes \(S_{n}\) a principal covering with group \(\{1,-1\}\) , this group acting by homotheties; the orbits are the pairs of diametrically opposite points.

\begin{enumerate}
\def\labelenumi{\arabic{enumi})}
\setcounter{enumi}{1}
\item
  Every covering of degree 2 has a unique structure of a principal covering with group Z/2Z.
\item
  Let X be a finite-dimensional locally separated differentiable manifold over R, and \(\widetilde{X}\) be the manifold of orientations of the tangent bundle of X (VAR, R, 10.2.4). The space \(\widetilde{X}\) , equipped with its canonical projection onto X, is a principal covering of X with group \(\{1,-1\}\) .
\end{enumerate}

COROLLARY 1. --- Let G be a discrete topological group acting continuously on the right on a topological space E. The following conditions are equivalent:

\begin{enumerate}
\def\labelenumi{(\roman{enumi})}
\item
  The group G acts properly and freely on E;
\item
  The space E/G is Hausdorff, and E is a principal covering with base E/G and group G.
\end{enumerate}

Moreover, under these conditions, the space E is Hausdorff.

Suppose condition (i) is satisfied. Then the spaces E and E/G are separated (TG, III, p. 29, prop. 3), and for every \(x \in \mathrm{E}\) there exists an open neighborhood U of \(x\) in E such that \(\mathrm{U} \cap (\mathrm{U} \cdot g) = \varnothing\) for every element \(g\) of G other than the identity element (TG, III, p. 32, prop. 8). By th. 1, condition (ii) is then satisfied.

The implication (ii)⇒(i) follows from Corollary 1 of I, p. 96.

COROLLARY 2. --- Let G be a topological group and let H be a discrete subgroup of G. Let H act on G by right translations. Then the canonical map from G into the space G/H of left cosets modulo H endows G with a structure of a principal covering of G/H with group H.

Let V be a neighborhood of the identity element \(e\) of G such that \(\mathrm{H} \cap \mathrm{V} = \{e\}\) . There exists an open neighborhood U of \(e\) in G such that \(\mathrm{U}^{-1} \cdot \mathrm{U} \subset \mathrm{V}\) (TG, III, p. 3) and, consequently, \(\mathrm{U} \cap \mathrm{U} \cdot h = \varnothing\) for all \(h \in \mathrm{H}\) , \(h \neq e\) . Corollary 2 of I, p. 100 therefore follows from Theorem 1.

Example 4. --- Equipped with the canonical map from R to T = R/Z (TG, V, p. 2), R is a principal covering of R/Z with group Z.

Remark 1. --- Let G be a separated topological group, K a compact subgroup of G, and H a discrete subgroup of G. The group G acts continuously and properly on the right on the space K\G of right cosets modulo K (TG, III, p. 30, corollary). Since the group H is closed in G, it also acts properly on K\G (TG, III, p. 27, example 1). If, moreover, one has \(H \cap gKg^{-1} = \{e\}\) for all \(g \in G\), the group H acts freely on K\G. The space K\G is then a principal covering with group H of K\G/H by Corollary 1.

COROLLARY 3. --- Let G and G' be topological groups, and let \(\varphi\colon G \to G'\) be a continuous, open, and surjective homomorphism. Suppose that the kernel H of \(\varphi\) is discrete. Then, for the action of H on G by right translations, \(\varphi\) makes G a principal covering of G' with group H.

If moreover the group G is connected, then H is contained in the center of G.

The homomorphism \(\varphi\) induces an isomorphism of topological groups from G/H onto G' (TG, III, p. 16, prop. 24), whence the first assertion by Corollary 2. Suppose the group G is connected. For every \(h \in \mathrm{H}\) , the continuous map from G into H defined by \(g \mapsto ghg^{-1}\) is constant, with value \(h\) . The group H is therefore contained in the center of G.

Remarks. --- 2) Let G and G' be topological groups, and let \(\varphi\colon\mathrm{G}\to\mathrm{G}'\) be a continuous and open homomorphism. If the group G' is connected, the homomorphism \(\varphi\) is surjective. Indeed, \(\varphi(\mathrm{G})\) is an open, hence closed, subgroup of \(\mathrm{G}'\), and consequently is equal to \(\mathrm{G}'\).

\begin{enumerate}
\def\labelenumi{\arabic{enumi})}
\setcounter{enumi}{2}
\tightlist
\item
  Let G be a locally compact group which is countable at infinity, and let G' be a separated topological group whose underlying space is a Baire space. Every continuous surjective homomorphism from G onto G' is open (TG, IX, p. 56, corollary and TG, III, p. 16, prop. 24).
\end{enumerate}

Examples. --- 5) For every integer \(n > 0\), let us denote by \(\mu_{n}\) the group of n-th roots of unity in C (A, V, p. 75). The map \(z \mapsto z^{n}\) makes \(C^{*}\) a principal covering of \(C^{*}\) with group \(\mu_{n}\), this group acting in \(C^{*}\) by multiplication.

\begin{enumerate}
\def\labelenumi{\arabic{enumi})}
\setcounter{enumi}{5}
\tightlist
\item
  Equipped with the map \(z \mapsto e^{2\pi iz}\) (TG, VIII, p. 8, remark), the space C is a principal covering of \(C^{*}\) with group Z. The map \(x \mapsto e^{2\pi ix} = \mathbf{e}(x)\) from R to \(S_{1}\) makes R a principal covering of \(S_{1}\) with group Z.
\end{enumerate}

